\chapter*{Related Work, and Where This Differs}
\label{ch:related}
\addcontentsline{toc}{chapter}{Related Work, and Where This Differs}
\markboth{Related Work}{Related Work}
\renewcommand{\thetable}{\Alph{table}}\setcounter{table}{0}
% ===================================================================

A related-work section normally goes at the back, where it can be
read by people who already agree with you.
This one is at the front because the book is long, and because a
reader is entitled to know before investing which existing
programme this most resembles and why it is not that one.

There is a line to carry through all of it.
Of every programme discussed below, this is the only one in which
the knower's epistemic limits are \emph{priced}.
Not bounded, not idealised away, not left to an implementation ---
priced, in a conserved resource, on a ledger the agent cannot
refill from outside.
Everything else i might have offered as the differentiator turns
out to be shared with somebody.
That one is not, and the rest of this chapter is partly an attempt
to check that claim against the strongest neighbours i can find.

For each neighbour i try to say four things: what it claims, what
this book takes from it, where they part company, and whether the
parting could be checked by anything.
The last is the one that makes this more than a bibliography with
opinions.

\section*{Theories of mind and consciousness}

\paragraph{Integrated Information Theory.}
Tononi's $\Phi$ is a graded measure of consciousness computed from
a system's causal structure \cite{tononi2008consciousness}, and
this book also grades rather than thresholds --- a point
Chapter~\ref{sec:hyper-discussion} used to make in a sentence and
which deserves more.
The parting is over what the grading is a measure \emph{of}.
$\Phi$ is intrinsic: it is a property of a system's own cause-effect
structure, computed by an analyst standing outside it.
The coordinate here is relational: it is a position in a lattice of
choice strength, and it is a fact about what an agent can do at a
cut with another agent on the far side.
The two come apart on a case that is easy to state.
A system with rich internal causal structure and no interaction
surface --- nothing it can meet, nothing it can be met by --- has,
on IIT's account, whatever $\Phi$ its structure gives it.
On this account it has no coordinate at all, because a coordinate
is a fact about apertures and it has none.
Whether that is a reductio of one view or the other is exactly the
question, and it is at least a question with an answer.
A second and less flattering difference: IIT has an experimental
programme and this does not.

\paragraph{The free energy principle.}
Friston's account \cite{friston2010free} is the closest neighbour
in the book, and the resemblance is real: both make prediction the
central activity of a persisting system, and both make the failure
of prediction the thing that threatens persistence.
The parting is that the free energy principle has no ledger.
Prediction error is minimised, but nothing is \emph{spent} to
minimise it, and there is no account that can run dry.
That difference is not cosmetic, and Chapter~\ref{ch:sci} shows
where it bites: once assays cost, being right stops being the same
thing as being fed, and a cruder theory that costs a fraction as
much can dominate the true one.
An agent minimising prediction error alone would not do that.
An agent on a budget provably will.

The honest concession is in the other direction.
A single scalar surplus is a thin theory of motivation next to an
affective landscape, and this book's $\varsigma$ is a budget line
rather than a feeling: one drive, no valence, no arousal, nothing
that is \emph{about} anything.
Chapter~\ref{ch:eng} is where that gap is partly closed --- token
colours, engines, and the virtual token are a theory of how
incommensurable drives become commensurable --- but it is not yet
a theory of affect and should not be read as one.

\paragraph{Global workspace.}
Baars's architecture \cite{baars1988} and Dehaene's neuronal
implementation \cite{dehaene2011} make consciousness a matter of
broadcast: content becomes conscious by winning access to a shared
channel that many specialised processes can read.
Chapter~\ref{ch:com} tells a broadcast-and-contention story of its
own, and the resemblance is close enough to be worth naming.
The parting is over where the architecture comes from.
Global workspace theory posits the workspace, because brains appear
to have one.
Here the composite is \emph{forced}: a single learner revising its
theory incrementally provably cannot leave the region it started
in, so the unit of learning has to become a population sharing a
namespace, and the shared namespace is the workspace arriving as a
consequence rather than as a design.

\paragraph{Attention schema.}
Graziano's proposal \cite{graziano2013} is that awareness is the
brain's model of its own attention --- a self-model, and a
deliberately imperfect one.
This is the neighbour whose absence i feel most.
The germ carried by a learner in Chapter~\ref{ch:sci} \emph{is} a
self-model, and the book uses it almost exclusively for
reproduction.
Worse, the machinery for the rest is already present:
Section~\ref{sec:kant-affordable}'s argument, applied reflexively,
gives an agent that cannot afford to resolve its own structure and
therefore reports a simplified one --- which is the introspective
illusion, derived rather than assumed, in about three sentences.
Those three sentences are not in this book.
They should be.

\paragraph{Chalmers, and the hard problem.}
The Pledge argues that the productive question is not what it is
like to be a bat but what it is like to be a computer program
\cite{nagel1974,chalmers1996}.
It should be said plainly what that does and does not accomplish.
It does not dissolve the hard problem.
It relocates it into a setting where the surrounding questions
become tractable --- what an agent can represent, what it can
afford to represent, what it must be committed to, what it cannot
distinguish and cannot tell that it cannot distinguish --- and
leaves the residue exactly where Chalmers left it.
Making a question tractable is worth doing.
It is not the same as answering it, and this book does rather more
of the first than the second.

\section*{Computational and physical foundations}

\paragraph{Wolfram's physics project and the ruliad.}
The nearest neighbour in ambition \cite{wolfram2020,wolfram2021}:
space, time, and a good deal of physics derived from rewriting
rather than presupposed.
Two partings, and both are structural.
First, a GSLT carries equations and a named site of interaction, so
the cut at which two terms meet is part of the presentation rather
than something that emerges from the rewriting; that is what makes
it possible to meter anything at all.
Second, and more sharply: this book fibres over a \emph{category}
of theories, in which no object is privileged.
The definite article in ``the ruliad'' is doing work that this
framework declines to do.
The rule space here is a variety of possible physics
(Part~\ref{part:physics}), and which one you get depends on which
further axioms hold --- a difference that shows up immediately in
what each programme thinks it is explaining when it recovers a
piece of physics.

\paragraph{Constructor theory.}
Deutsch and Marletto's programme \cite{deutsch2013,marletto2015}
is the nearest neighbour in \emph{method}, and the resemblance is
striking once seen.
Both replace statements about how a system evolves with statements
about which transformations are possible.
Both take that reformulation to be more fundamental than the
dynamical laws it subsumes.
And the parting is unusually clean.
Constructor theory's possible/impossible is Boolean.
Here it is priced.
That is not an analogy: it is exactly the distinction
Chapter~\ref{sec:weighted} already draws between decorating a
theory in the Boolean semiring and decorating it in the tropical
one.
Constructor theory is, on this reading, the Boolean fibre of a
construction this book runs over a family of semirings, and
possible-at-a-cost is the tropical one.
If that identification survives being made carefully, it is the
most useful single sentence in this chapter, and i do not think it
has been made before.

\paragraph{It from bit, and the computational universe.}
Wheeler \cite{wheeler1990} and Lloyd \cite{lloyd2006} are
ancestors rather than rivals, and the debt is general.
What this book adds to the slogan is a specific answer to
\emph{which} computation, argued for in
Section~\ref{sec:pledge-viii} rather than assumed: not a function
machine, because the world is interactive, and not merely an
interactive one, because an agent has to be able to represent its
own hypotheses.

\paragraph{Applied category theory.}
The categorical machinery here is standard and the debts are to a
tradition rather than to a result: Spivak and Fong's presentation
of the field \cite{fongspivak2019}, Coecke and Kissinger's
diagrammatic method \cite{coeckekissinger2017}.
One honest gap belongs in this paragraph.
The cost monad and the history monad are applied together
throughout the book, and no distributive law between them is ever
stated.
Remark~\ref{rmk:two-erasures}, on the two erasures, is very close to
being the observation that they interchange only laxly, and it
stops short of saying so.

\paragraph{The technical lineage.}
Milner's process calculi \cite{milnerCC}, Girard's linear logic,
Abramsky's proofs-as-processes \cite{abramskyPaP}, Turi and
Plotkin's bialgebraic semantics \cite{turiplotkin}, Turing's oracle
hierarchy \cite{turing1939}, and the Weihrauch degrees
\cite{weihrauch} are the actual apparatus, cited throughout and not
argued with.
The one place the book takes a position against the tradition is
Chapter~\ref{ch:chs}, which claims that linear logic types the
communication skeleton of a computation and is silent by
construction about how its nondeterminism resolves --- so that a
logic of choice is the complement of linear logic rather than a
refinement of it.

\section*{Learning, agency, and the symbol problem}

\paragraph{Solomonoff induction and AIXI.}
Solomonoff's prior \cite{solomonoff1964} and Hutter's AIXI
\cite{hutter2005} are the canonical account of optimal prediction
and optimal action given unbounded computation.
The mortal scientist is the same problem with a budget and a body,
and the relationship is sharper than a slogan.
Proposition~\ref{sci:prop:limit} says the complete theory of an
environment sits at the top of the lattice as a limit no learner
with a finite token supply can install --- which is to say that
AIXI's optimality is precisely the thing this framework proves to
be unpurchasable.
The two accounts are not competitors so much as the two ends of one
axis, and the interesting content is entirely in what happens
between them.

\paragraph{Computational mechanics.}
Crutchfield and Shalizi's $\epsilon$-machines
\cite{shalizi2001} are load-bearing in the Prestige rather than
merely adjacent.
Causal states are the coarsest coarse-graining of pasts that
retains full predictive power, and the proposal of
Chapter~\ref{ch:notation} is that they are the natural first
candidate for which distinctions in a host world deserve builtin
names.
The taxonomy of $\epsilon$-machines --- finite, countable,
uncountable, fractal --- turns out to be a density hierarchy in
Goodman's sense, which converts a question about notational
adequacy into a question about whether a budget-quotiented
$\epsilon$-machine is finite.
Both sides of that equivalence are computable.
Neither has been computed.

\paragraph{Symbol grounding and emergent communication.}
Harnad's problem \cite{harnad1990}, Goodman's theory of
notationality \cite{goodman1968}, Steels's language games
\cite{steels2015}, and Taniguchi's collective predictive coding
\cite{taniguchi2024} are treated at length in
Chapters~\ref{ch:trick} and~\ref{ch:notation} and are only noted
here.
The short version: the correspondence between a language and a
world is a fact about the language's users, the emergent
communication literature supplies the population that a
single-agent account structurally cannot, and neither solves the
problem this book concedes it has not solved.

\paragraph{Bisimulation metrics.}
The reinforcement learning literature has independently hit the
rigidity of exact bisimulation and answered with quantitative
relaxations \cite{ferns2004,gelada2019,zhang2021,castro2021}.
Observation~\ref{obs:metric} argues this is the right repair in
that setting and the wrong one here, because a metric dissolves
exactly the discreteness a notation requires.
This is a genuine disagreement with active work and i would rather
state it as one than leave it in a footnote.

\paragraph{Hyperon and MeTTa.}
Goertzel's programme \cite{goertzel2021} is the closest thing to a
sibling project, and its near-absence from this book is the most
conspicuous omission in this chapter.
Both take a rewriting substrate to be the right foundation for
general intelligence; both make reflection central; and MeTTaIL,
the language-definition platform much of the machinery here was
implemented against, sits directly between them.
The comparison i owe and have not made is between MeTTa's
approach to self-modification and this book's --- reproduction as
recombination on quoted code --- and between Hyperon's
attention allocation and the metering here.
That is a chapter, not a paragraph, and it is not written.

\section*{Life and origins}

\paragraph{Assembly theory.}
Cronin and Walker's programme \cite{cronin2021,walker2023} is the
largest single engagement in the book.
What is taken: the assembly index as a measure of causal depth ---
the minimum number of steps by which an object can be built with
reuse --- and the claim that depth combined with abundance is a
biosignature.
What is offered back: answers to two questions the theory leaves
open, namely what counts as a copy and how to specify the scope
within which copies are counted, both of which
Chapter~\ref{sec:scope-manufacture} treats as fixed-point problems
with computable answers.
Where it dissents: Proposition~\ref{scm:prop:geom} says that a
coarse instrument does not merely mismeasure copy number, it
\emph{manufactures} copies, at a rate exponential in the coarseness,
and the manufactured copies are biased toward the biosignature the
measurement is looking for.
That is a criticism of the measurement protocol rather than of the
theory, and it should be read as such.
It is also, so far as i know, the only claim in this book that a
laboratory could check next year.

\paragraph{Smith and Morowitz.}
The phase-transition account of the origin of life
\cite{SmithMorowitz} supplies the picture of a thermodynamically
driven whole within which chemistry is channelled.
Chapter~\ref{ch:agy} reaches the conclusion that Smith's
thermodynamic macro-component and the determinacy islands of this
framework \emph{cross} rather than nest --- and an earlier reading
of my own, which had them nesting, is retained in a remark
explaining why it does not survive.

\paragraph{Autopoiesis and $(M,R)$ systems.}
Maturana and Varela \cite{maturanavarela1980} and Rosen
\cite{rosen1991} are the prior attempts at organisational closure,
and Rosen is the one that should have been in this book from the
start.
``Closed to efficient causation'' --- the requirement that every
function of an organism be produced by the organism --- is close
kin to the replication-as-a-fixed-point construction of
Chapter~\ref{sec:replication}, close enough that the difference is
worth stating.
Rosen's closure is a condition on a category of mappings, and his
conclusion is that such systems are not simulable by machines.
The closure here is a fixed point of a reflective rewrite theory,
and the same structure is exhibited constructively.
Whichever of us is wrong, we are wrong about something specific,
which is more than most pairs of positions in this area manage.

\section*{The table}

Table~\ref{tab:related} scores the programmes on five questions.
The last column is where the argument of this chapter lives, and it
is a claim rather than a description --- which is why the column
exists and why it is last.

\begin{table}[htbp]
\centering
\small
\renewcommand{\arraystretch}{1.35}
\begin{tabular}{@{}p{2.1cm} p{2.3cm} p{2.2cm} p{1.5cm} p{2.4cm} p{1.6cm}@{}}
\toprule
\textbf{Programme} & \textbf{Theory of} & \textbf{Primitive} &
\textbf{Observer inside?} & \textbf{What would falsify it} &
\textbf{Anything priced?} \\
\midrule
IIT & consciousness & cause--effect structure & no &
a high-$\Phi$ system reporting nothing & \textsc{no} \\
Free energy & persistence & prediction error & yes &
persistence without error minimisation & \textsc{no} \\
Global workspace & access consciousness & broadcast &
partly & conscious report without broadcast & \textsc{no} \\
Constructor theory & physical law & possible / impossible
transformations & no & a task both possible and impossible &
\textsc{boolean only} \\
Ruliad & physics & rewriting & yes & a physics not in the rule
space & \textsc{no} \\
AIXI & optimal agency & algorithmic probability & no &
--- (optimal by construction) & \textsc{no} \\
Assembly theory & life detection & causal depth $\times$
abundance & no & high assembly index abiotically &
\textsc{depth, not access} \\
\midrule
\textbf{This book} & all of the above, at a cost &
metered interaction at a cut & yes &
finite contention; the copy-manufacture bound &
\textsc{yes} \\
\bottomrule
\end{tabular}
\caption{Five questions, scored. The final column is the claim this
chapter is making: the framework developed here is the only one on
the list in which what an agent can know is limited by what it can
afford rather than by a prohibition. Constructor theory scores
partially because possible/impossible is a pricing with two prices.
Assembly theory scores partially because it prices \emph{construction}
and says nothing about the cost of \emph{access}.}
\label{tab:related}
\end{table}

Two rows deserve comment because they are the ones i expect
argument about.

Constructor theory is marked ``Boolean only'' rather than
``no'', and that is deliberate: a two-valued possibility measure is
still a measure, and on the reading proposed above it is precisely
the Boolean fibre of the semiring-parametric construction of
Chapter~\ref{sec:weighted}.
Whether that is a generalisation of constructor theory or a
misreading of it is a fair question and i do not think the answer
is obvious.

AIXI's falsification cell is a dash, and this is not a cheap shot.
AIXI is optimal by construction relative to its own prior, so there
is no observation that could refute it; what can be refuted is the
claim that it is a useful model of anything realisable, and
Proposition~\ref{sci:prop:limit} is one form of that refutation.
A framework that cannot be wrong is doing something other than what
this book is trying to do, and the dash records that rather than
scoring it.

\section*{Two concessions}

\paragraph{Where this is a special case rather than a rival.}
The account of interaction, the modal logic, and the categorical
apparatus are all standard, and nothing in
Part~\ref{part:machinery} is offered as a departure from the
process-calculus tradition.
The novelty, if there is any, is in what happens when a meter is
installed and an agent is made of the same material as its
environment.
A reader who wants to describe the whole of the Machinery as
``known, assembled unusually'' will get no argument here.

\paragraph{Where a neighbour is ahead.}
Two of these programmes have empirical arms and this one does not.
The free energy principle has a substantial experimental
literature.
IIT has a perturbational complexity index and clinical data behind
it.
Assembly theory has a mass spectrometer and a published protocol.
This book has falsification conditions --- finite contention in the
massive-population argument, and the copy-manufacture bound of
Chapter~\ref{sec:scope-manufacture} --- and they are Fermi
estimates rather than experiments.
The second of the two is the closer to being run, and it is the
thing i would most like someone to take away from this chapter.

%% restore the book's numbering for the numbered chapters that follow
\renewcommand{\thetable}{\thechapter.\arabic{table}}\setcounter{table}{0}
